%!TEX root = ../memoria.tex
% =============================================================================
%  CAPÍTULO 1 — INTRODUCCIÓN
% =============================================================================

Este primer capítulo presenta el contexto y la motivación del proyecto, define
los objetivos que se persiguen y describe la organización del resto de la
memoria. El punto de partida es la brecha digital que afecta a las personas mayores. Este problema social guia cada una de las decisiones de diseño, implementación y evaluación desarrolladas en
los capítulos siguientes.

\section{Motivación\label{SEC:MOTIVACION}}

Las redes sociales, la mensajería instantánea y las
videollamadas han convertido el contacto
con familiares y amigos en una cuestión de segundos. Sin embargo, este avance
también ha expuesto una brecha entre generaciones. Las personas mayores
representan un gran desencuentro entre
tecnología y usuario, y es precisamente a ellas a quienes se dirige este proyecto.

El fenómeno, conocido como brecha digital, está ampliamente documentado en España. Los datos del Instituto Nacional de Estadística revelan que una
proporción muy significativa de la población mayor de 65 años reside en hogares
unipersonales, mientras que los estudios sobre soledad también hablan de grandes números de personas en este rango de edad que han experimentado
soledad no deseada \cite{ine_hogares,rsi_soledad}. Las estudios sobre las consecuencias de la soledad destacan: deterioro del descanso, aceleración del declive
cognitivo y mayor predisposición a enfermedades neurodegenerativas
\cite{soledad_cerebro}. A ello se suman la distancia respecto a los familiares
más cercanos y las limitaciones físicas asociadas al envejecimiento.

La tecnología tiene un potencial evidente para reducir esa distancia, pero solo
si se diseña con ese pensamiento. Las plataformas de uso cotidiano,
como WhatsApp o las grandes redes sociales, no fueron diseñadas para las personas
mayores: incorporan funcionalidades nuevas de forma
constante, modifican su interfaz con frecuencia, mezclan acciones distintas en una
misma pantalla y emplean iconografía dificl de entender si no estas acostumbrado a ella. Por un lado, muchas personas mayores que se beneficiarían
de estar conectadas a la tecnología terminan dependiendo de un familiar para realizar acciones
básicas. Por otro, otras renuncian directamente al uso de la tecnología.

\emph{ConectaMayor} parte de la idea de que es posible construir una herramienta
digital sencilla, útil y privada, capaz de facilitar el contacto cotidiano entre
las personas mayores y su entorno familiar. La aplicación no tiene el objetivo de competir con el resto de plataformas. Pero ofrece, en su lugar, un espacio digital sencillo
en el que la persona mayor puede gestionar su agenda, consultar sus contactos,
intercambiar mensajes con un familiar o ver las fotografías compartidas de la familia.

\section{Objetivos\label{SEC:OBJETIVOS}}

El objetivo general de este Trabajo Fin de Grado es \textbf{diseñar, desarrollar
y evaluar una aplicación web accesible que permita a las personas mayores
comunicarse con sus familiares en un espacio digital privado, sencillo y seguro},
aplicando los principios del diseño centrado en el usuario en todas las fases del
proyecto.

A partir de este objetivo general, se han establecido los siguientes objetivos
específicos:

\begin{enumerate}
  \item Analizar el contexto del problema y revisar las soluciones existentes desde el punto de vista de la accesibilidad para
  personas mayores.

  \item Diseñar una arquitectura organizada en módulos funcionales: agenda, contactos, mensajería, galería y vinculación familiar.Con un sistema de roles que adapte los permisos y la interfaz al perfil de cada
  usuario.

  \item Implementar una interfaz diferenciada para el rol de persona mayor, con
  tipografía ampliada, áreas de interacción grandes, lenguaje sencillo y flujos
  de navegación reducidos al mínimo, manteniendo en paralelo una interfaz
  completa para los familiares.

  \item Cumplir las pautas de accesibilidad WCAG 2.1 en su nivel AA, aplicándolas
  a colores, tamaños, contrastes, indicadores de foco y mecanismos de interacción.

  \item Garantizar la privacidad mediante un modelo de acceso restringido en el
  que los datos de cada familia permanezcan aislados del resto.

  \item Evaluar la usabilidad con personas mayores siguiendo la metodología de la disciplina de Interacción Persona-Ordenador (IPO), e incorporar a la aplicación las mejoras detectadas durante las sesiones.

  \item Documentar el proceso, las decisiones técnicas y los resultados de la
  evaluación.
\end{enumerate}

\section{Estructura del documento\label{SEC:ESTRUCTURA}}

La memoria se organiza en cinco capítulos y dos apéndices.

El \textbf{Capítulo~\ref{CAP:INTRODUCCION}} introduce el proyecto. En él se
presenta la motivación del trabajo, centrada en la brecha digital que afecta a
las personas mayores y en la necesidad de diseñar herramientas tecnológicas que ayuden a su comunicación con el entorno familiar. Además, se
definen el objetivo general y los objetivos específicos del proyecto, y se
describe la estructura seguida en la memoria.

El \textbf{Capítulo~\ref{CAP:ESTADOARTE}} presenta el estado del arte. Se
analizan las principales aplicaciones existentes que ofrecen soluciones
relacionadas con el problema abordado y se discuten sus limitaciones relaccionadas con la accesibilidad para personas mayores. Este análisis sitúa
\emph{ConectaMayor} en la actualidad y justifica la idea del
proyecto.

El \textbf{Capítulo~\ref{CAP:DISENO}} aborda el diseño de la aplicación. Se
exponen los requisitos funcionales y no funcionales agrupados por subsistemas,
se describe el modelo de datos, se presentan los prototipos de las pantallas
principales y se justifican las decisiones de arquitectura más importantes.

El \textbf{Capítulo~\ref{CAP:IMPLEMENTACION}} describe la implementación. Se
detallan las tecnologías empleadas, la arquitectura elegida, se presentan los
módulos desarrollados mediante capturas reales y fragmentos de código, y se
incluye, como sección final, la evaluación de usabilidad realizada con personas
mayores siguiendo la metodología IPO.

El \textbf{Capítulo~\ref{CAP:CONCLUSIONES}} presenta las conclusiones. Se valora
el grado de cumplimiento de los objetivos, se discute las limitaciones
encontradas y se proponen ideas y mejoras para el futuro.

El Apéndice~\ref{APP:AIPO} reúne las tablas detalladas de las pruebas de
usabilidad, cuyo análisis se ha integrado en el cuerpo del
Capítulo~\ref{CAP:IMPLEMENTACION}. El Apéndice~\ref{APP:ENTORNO} recoge el
entorno de ejecución y las instrucciones de instalación necesarias para
reproducir el proyecto.

El código fuente completo de la aplicación está disponible públicamente en el
repositorio \url{https://github.com/yagogl/conectamayor-TFG}, donde se incluyen
las instrucciones de instalación y ejecución. Los detalles del entorno y las
dependencias se recogen en el Apéndice~\ref{APP:ENTORNO}.